\begin{figure}[ht]
\centering
\begin{tikzpicture}[scale=0.85,
  gridln/.style = {draw=enggray!45, thin},
  diag/.style   = {draw=engred, thick, dashed},
  goodpath/.style = {draw=engblue, very thick, ->, >=stealth},
  badpath/.style  = {draw=engorange, very thick, ->, >=stealth},
  dot/.style    = {circle, fill=engblack, inner sep=1.1pt},
]
% 4x4 lattice grid.
\foreach \x in {0,1,2,3,4} {
  \draw[gridln] (\x,0) -- (\x,4);
  \draw[gridln] (0,\x) -- (4,\x);
}
% The forbidden diagonal boundary y = x + 1 (a path staying weakly
% below the main diagonal must never cross above this line).
\draw[diag] (0,1) -- (3,4);
\node[engred, anchor=west, font=\scriptsize] at (3.05,3.9)
  {barrier $y = x + 1$};

% A valid monotone lattice path (never rising above the diagonal):
% steps R,R,U,R,U,U,R,U in an x-then-y walk from (0,0) to (4,4).
\draw[goodpath] (0,0) -- (1,0) -- (2,0) -- (2,1) -- (3,1)
  -- (3,2) -- (3,3) -- (4,3) -- (4,4);
\node[engblue, anchor=north west, font=\scriptsize] at (0.05,-0.15)
  {valid Dyck path};

% A path that touches the barrier (would be rejected).
\draw[badpath] (0,0) -- (0,1) -- (1,1) -- (1,2);
\node[engorange, anchor=south east, font=\scriptsize] at (1.0,2.05)
  {rejected: crosses barrier};

% Corner dots.
\node[dot] at (0,0) {}; \node[dot] at (4,4) {};
\node[anchor=north east, font=\scriptsize] at (0,0) {$(0,0)$};
\node[anchor=south west, font=\scriptsize] at (4,4) {$(4,4)$};
\end{tikzpicture}
\caption{The Catalan number $C_{n}$ counts the monotone lattice
paths from $(0,0)$ to $(n,n)$ that never rise above the main
diagonal, the blue path being one of the $C_{4}=14$ valid paths on
the $4\times 4$ grid. The reflection principle counts the paths that
do cross the barrier $y=x+1$ (orange) by reflecting the offending
prefix, turning a forbidden endpoint into $(n-1,n+1)$; subtracting
those from all $\binom{2n}{n}$ monotone paths leaves
$\binom{2n}{n}-\binom{2n}{n+1}=\frac{1}{n+1}\binom{2n}{n}$. The same
count enumerates balanced bracket strings and binary-tree shapes.}
\label{fig:vol02:ch17:catalan-lattice}
\end{figure}
